\section{The actual complex exponential and cubical Hermitian frame}
\label{sec:faithful-complex-realisation}
We follow the source construction with the same algebraic abelian
variety, exponential, lattice and rigidified line. These first
steps construct the exponential and its cubical normal frame.
The last subsection supplies positivity, the perfect pairing and the
all-level section correspondence from its five explicit algebraic laws.
The relative-family and reverse algebraization comparisons remain separate.


\subsection{The actual complex exponential and its normalized tangent map}
\label{subsec:ftb-actual-complex-exponential}

This subsection follows the frozen source construction, before constructing the
Appell--Humbert frame.  Fix a polarised abelian variety
\(P\in\mathrm{Out.PolAV}(\operatorname{Spec}\mathbb C)\), of dimension \(g\).
We use the source's ample line bundle, affine charts and regular-function
evaluations.  The multiplication law used here is
\(\mathrm{Out.MulNSt}\): in particular, for each \(n\geq1\), multiplication
by \(n\) is a finite flat surjective morphism.  This is a previously supplied
algebraic law, not a consequence of the analytic construction below.

Write \(P(\mathbb C)\) additively.  Give it precisely the topology generated by
\[
 \{x:x\in U(\mathbb C),\ f(x)\in O\},
 \qquad U\subset P\text{ open},\quad f\in\Gamma(U,\mathcal O_P),
 \quad O\subset\mathbb C\text{ open}.                         \tag{CE1}
\]
For a map \(a:W\to P(\mathbb C)\), the source's holomorphy condition means that
\(a^{-1}U\) is open in the open set \(W\), and every
\(f\circ a\) is holomorphic on \(a^{-1}U\).  All occurrences of
holomorphic maps to algebraic varieties below have this meaning.

\begin{srlemma}[The affine and smooth charts used by the construction]
\label{lem:ftb-complex-algebraic-charts}
For an affine open \(U\subset X\) of a complex scheme, a
\(\mathbb C\)-algebra homomorphism
\(\Gamma(U,\mathcal O_X)\to\mathbb C\) determines a unique complex point of
\(U\), with these regular-function values.  If \(X\to\operatorname{Spec}\mathbb C\)
is smooth of relative dimension \(g\), then at any \(x\in X(\mathbb C)\)
there are an affine open \(U\ni x\), functions \(t_1,\ldots,t_g\), an open
neighbourhood \(W\subset\mathbb C^g\) of zero, and a map
\(\sigma:W\to X(\mathbb C)\), such that
\[
 \sigma(0)=x,\qquad t_i(\sigma(w))=w_i,                       \tag{CE2}
\]
\(\sigma\) is holomorphic, and \(\sigma(W)\) is open for (CE1).
Every derivation at \(x\) that kills the constants and all the \(t_i\)
is zero.
\end{srlemma}

\begin{proof}
The affine point is \(\operatorname{Spec}(e)\) followed by the affine-open
map.  The identities between pullback of sections and evaluation show both
its prescribed values and uniqueness; no comparison of point sets is
assumed.

Choose a standard smooth affine presentation near \(x\).  Its etale
coordinates can be translated by their values at \(x\), so that they
vanish there.  The analytic inverse used for this etale map is obtained
as follows.  Present the etale algebra over
\(\mathbb C[T_1,\ldots,T_g]\) by equally many extra variables and
relations, with invertible Jacobian in those extra variables.  If the
extra variables are \(Y\) and the relations are \(R\), apply the analytic
inverse-function theorem to
\[
             (w,y)\longmapsto (w,R(w,y)).
\]
The derivative is injective: its first component forces the variation
of \(w\) to be zero, and the invertible relation Jacobian then forces the
variation of \(y\) to be zero.  Source and target have equal finite
dimension.  Restrict the analytic inverse to \((w,0)\).
Its values satisfy every relation and hence give algebra homomorphisms
\(\sigma_w\) from the etale algebra to \(\mathbb C\).
For every algebra element, its value under \(\sigma_w\) is analytic,
because the element is represented by a polynomial in the chosen
variables.  The inverse identities give (CE2).

Here openness in (CE1) is also verified.  A finite presentation embeds
the set of algebra homomorphisms into a finite-dimensional coordinate
space.  Shrink the chart using finitely many tests
\[
       |\sigma_w(f_j)-e(f_j)|<\varepsilon .
\]
The analytic inverse contains every algebra homomorphism satisfying
these tests.  Thus its image contains the corresponding basic
neighbourhood, and after this restriction its image is such an open
chart.  A regular function on a different open set is locally a
quotient \(a/h^m\) on a basic open of \(U\).  Its value is the holomorphic
quotient \(\sigma_w(a)/\sigma_w(h)^m\), where the denominator is
nonzero.  This proves the stated holomorphy on every algebraic open.

Finally, a derivation killing the \(t_i\) kills their whole polynomial
algebra.  The etale algebra is formally unramified over that polynomial
algebra.  Apply uniqueness of lifts to the two maps into
\(\mathbb C[\epsilon]/(\epsilon^2)\), namely \(e+\epsilon D\) and \(e\).
They agree, so \(D=0\).
\end{proof}

Two further topological facts are established using these same values.
First, products of continuous, respectively holomorphic, families are
continuous, respectively holomorphic.  On affine product charts the
coordinate algebra is generated by the images of the two factor
algebras.  Addition and multiplication of their values prove the
assertion there.  Restriction to arbitrary opens is proved by
localization, using \(a/h^m\).  Applying this to the multiplication and
inverse morphisms makes \(P(\mathbb C)\) a topological group.

Second, this topological space is compact.  Here is the source's
projective-chart proof.  Choose finitely many sections of an ample
power that cover \(P\) by their nonvanishing loci, such that the ratios
of these sections generate the affine coordinate rings.  In each
chart, the homomorphisms for which all these ratios have modulus at
most one form a compact set: they are the closed zero locus of the
relations inside a finite closed unit polydisc.  At every point choose
a section of maximal modulus among the chosen sections, in any local
frame.  Its nonvanishing chart has every ratio of modulus at most one.
Consequently these finitely many compact sets cover \(P(\mathbb C)\).
The affine point correspondence in
Lemma~\ref{lem:ftb-complex-algebraic-charts} identifies their values with
the actual points throughout this argument.

\begin{srlemma}[The source's halving logarithm]
\label{lem:ftb-halving-logarithm}
Let \(V\) be a finite-dimensional complex normed vector space.
Suppose \(m\) is an analytic commutative local group law at zero, with
unit zero and associative wherever the displayed compositions are
defined near zero.  There is an analytic local change of coordinates
\(\mathcal L\), with
\[
 \mathcal L(0)=0,\quad d\mathcal L_0=1,\qquad
 \mathcal L(m(x,y))=\mathcal L(x)+\mathcal L(y)                 \tag{CE3}
\]
for sufficiently small \(x,y\).
\end{srlemma}

\begin{proof}
The unit identities and commutativity give
\[
          m(x,y)=x+y+O(\|(x,y)\|^2).
\]
The derivative of \(d(x)=m(x,x)\) is \(2\), so the analytic
inverse-function theorem gives a local inverse \(\varepsilon\), with
\[
       \varepsilon(x)=\tfrac12x+O(\|x\|^2).
\]
Associativity and commutativity give
\(d(m(x,y))=m(d(x),d(y))\).  Applying the local inverse, after shrinking
all neighbourhoods, gives
\[
       \varepsilon(m(x,y))=m(\varepsilon(x),\varepsilon(y)).   \tag{CE4}
\]
Choose \(r>0\) and \(C\geq0\) such that, for \(\|x\|<r\),
\[
 \|\varepsilon(x)\|\leq\tfrac58\|x\|,\qquad
 \|\varepsilon(x)-\tfrac12x\|\leq C\|x\|^2.
\]
All iterates remain in this ball.  Put
\(\mathcal L_j(x)=2^j\varepsilon^{\,j}(x)\).  Direct subtraction gives
the exact estimates used in the source:
\[
 \|\mathcal L_{j+1}(x)-\mathcal L_j(x)\|
       \leq2C\|x\|^2(25/32)^j,\qquad
 \|\mathcal L(x)-\mathcal L_j(x)\|
       \leq\tfrac{64}{7}C\|x\|^2(25/32)^j,                  \tag{CE5}
\]
where \(\mathcal L=\lim_j\mathcal L_j\).
The second estimate follows by summing the geometric tail.
Convergence is uniform on a smaller ball; its terms are holomorphic.
The locally uniform limit theorem therefore makes \(\mathcal L\)
holomorphic there.  Taking \(j=0\) in (CE5) gives
\(\mathcal L(x)-x=O(\|x\|^2)\), so \(d\mathcal L_0=1\).

Iterate (CE4).  If
\(\|m(x,y)-x-y\|\leq D\|(x,y)\|^2\) near zero, the difference
\[
 2^j m(\varepsilon^{\,j}(x),\varepsilon^{\,j}(y))
 -2^j\varepsilon^{\,j}(x)-2^j\varepsilon^{\,j}(y)
\]
has norm at most
\(D\|(x,y)\|^2(25/32)^j\).  It tends to zero, proving (CE3).
Another application of the inverse-function theorem gives an analytic
local inverse to \(\mathcal L\).
\end{proof}

\begin{srproposition}[The actual exponential, lattice and tangent comparison]
\label{prop:ftb-actual-complex-exponential}
There are an additive surjection
\[
                 \exp_P:\mathbb C^g\longrightarrow P(\mathbb C)
\]
and a discrete integral lattice
\(\Lambda=\ker(\exp_P)\), whose real span is \(\mathbb C^g\).
The map \(\exp_P\) is holomorphic in the regular-function sense above.
One may choose the centred coordinates \(t_i\) at the origin so that
\[
             d_0(t_1\circ\exp_P,\ldots,t_g\circ\exp_P)=1.    \tag{CE6}
\]
For each \(w\in\mathbb C^g\) there is a unique algebraic tangent vector
\(v_w:\operatorname{Spec}\mathbb C[\epsilon]/(\epsilon^2)\to P\)
at the origin such that, for every regular function defined there,
\[
          f(v_w)=f(0)+\epsilon\,d_0(f\circ\exp_P)(w).         \tag{CE7}
\]
Every algebraic tangent vector occurs uniquely in this way.
Every algebraic group endomorphism \(f:P\to P\) admits a complex
linear lift \(F\) with
\[
                         \exp_P(Fz)=f(\exp_P z).            \tag{CE8}
\]
Moreover, for every \(n\geq0\) and every regular-value-holomorphic
map \(a:\mathbb C^n\to P(\mathbb C)\) with \(a(0)=0\), there is a
map \(G:\mathbb C^n\to\mathbb C^g\), analytic at zero, with
\(G(0)=0\) and \(\exp_P(Gw)=a(w)\) for all sufficiently small \(w\).
\end{srproposition}

\begin{proof}
Use Lemma~\ref{lem:ftb-complex-algebraic-charts} at the origin.
The local operation in its coordinates is analytic by the product
calculation above.  Compose the inverse of the logarithm in
Lemma~\ref{lem:ftb-halving-logarithm} with the inverse coordinate chart.
This is a local map \(e_0\) with
\(e_0(v+w)=e_0(v)+e_0(w)\) near \((0,0)\).
It has a unique global additive extension: for \(j\) sufficiently
large set
\[
                     e(v)=2^j e_0(2^{-j}v).                \tag{CE9}
\]
Repeated local doubling shows that the right side is independent of
\(j\).  For additivity, choose one \(j\) placing \(2^{-j}v\),
\(2^{-j}w\), and \(2^{-j}(v+w)\) in the common small neighbourhood.
Uniqueness follows by applying the same rescaling to any two additive
maps that agree near zero.  Set \(\exp_P=e\).

Near zero, \(e\) is the local analytic coordinate inverse just
constructed.  Hence it is a local homeomorphism there, and translation
makes it a local homeomorphism everywhere.  Its image is an open
subgroup.  The quotient by this image is a finite group, by compactness.
Multiplication by every positive integer is surjective on \(P(\mathbb C)\):
the finite surjective algebraic multiplication morphism has a closed
point in the fibre over each complex point, and its residue field is
\(\mathbb C\).  A finite divisible abelian group is zero (multiply by
its order).  Thus \(e\) is surjective.  The local homeomorphism at zero
also makes its kernel discrete.

For completeness, the source proves full real span directly.
The image is compact, so for some \(R>0\) every point in it has a
representative \(b\) with \(\|b\|<R\): the open sets \(e(B(0,n))\)
give an increasing cover, and a finite subcover suffices.
For any \(v\) and any large positive real \(a\), choose such \(b\)
with \(e(av)=e(b)\).  Then \(av-b\in\Lambda\), and
\[
             a^{-1}(av-b)\in\operatorname{span}_{\mathbb R}\Lambda,
             \qquad \|v-a^{-1}(av-b)\|<R/a .
\]
A real linear subspace in finite dimension is closed, so the span is
all of \(\mathbb C^g\).

Holomorphy at every \(z_0\) follows from
\(e(z)=e(z_0)+e(z-z_0)\) and holomorphy near zero; the actual product
and regular-function localization calculations above apply.
The inverse logarithm has derivative the identity, giving (CE6).

For the tangent assertion, define on an affine neighbourhood \(U\)
of zero
\[
 J_w(a)=a(0)+\epsilon\,d_0(a\circ e)(w).
\]
The product rule makes \(J_w\) a \(\mathbb C\)-algebra homomorphism;
the affine point correspondence over the dual numbers gives \(v_w\).
For an arbitrary tangent \(v\), take
\(w_i\) to be the \(\epsilon\)-coefficient of \(t_i(v)\).
The difference of the two corresponding derivations kills all
constants and the \(t_i\), so it is zero by the last assertion of
Lemma~\ref{lem:ftb-complex-algebraic-charts}.  This proves existence
and uniqueness on \(U\).  Localization to basic opens and restriction
of the evaluation homomorphisms proves (CE7) for every regular
function defined at the origin.

For the general local lifting assertion, write \(\varphi\) for the
origin coordinate chart used above.  The map \(\varphi\circ a\) is
analytic at zero by the regular-value definition and has value zero.
Put \(G=\mathcal L\circ\varphi\circ a\) near zero, and extend it
arbitrarily outside this neighbourhood.  The local inverse identities
give \(G(0)=0\) and \(e\circ G=a\) near zero.
Apply this assertion to \(a=f\circ e\) to obtain a holomorphic
\(G:\mathbb C^g\to\mathbb C^g\) near zero, with \(G(0)=0\) and
\(e(Gz)=f(ez)\).
The difference \(G(2z)-2G(z)\) belongs to the discrete kernel and
tends to zero.  It is therefore zero for all sufficiently small \(z\).
For a fixed \(z\), the sequence \(2^jG(2^{-j}z)\) is eventually
constant, while differentiability makes its limit \(dG_0(z)\).
Applying \(e\) and repeatedly doubling the identity
\(e(G(2^{-j}z))=f(e(2^{-j}z))\) proves (CE8) with \(F=dG_0\).
\end{proof}

\paragraph{Exact source scope and dependency ledger.}
The output here is the exponential part of
\srcid{AbelianExp.exists\_holUnif}, including its extra tangent-coordinate
and local lifting assertions; it initially uses the trivial line bundle
with constant frame \(1\).  It does not yet give the frame of the
polarising bundle, its positive Hermitian form, the Weil pairing formula,
or section--theta correspondence.  Those remain the subsequent,
separate steps of \srcid{thetaUnifAV\_of\_laws}.

All paths in this ledger are under
\srcid{src/Development/AbelianVarietyTheory/Proofs/} at tag
\srcid{alpoge-source-7cbf7a0c}.
The affine point proof is
\srcid{ToricGeometry/ExistsUniquePointOfRingHom.lean}
(blob \srcid{b4794489f6856874c710b89c8f1d3f65f8437af9}).
The analytic etale inverse is
\srcid{HomogeneousIdeals/ExistsOpenPartialHomeomorphAnalytic.lean}
(\srcid{3574b51a9512022e9702954181112212d94db05f});
its finite-test form is
\srcid{ToricGeometry/ExistsTestFunctionsSubsetOfMemNhds.lean}
(\srcid{5ba21a914397f864d2e7b2660a78dd12a9d8e243}), and the smooth
algebraic charts and derivation uniqueness are
\srcid{ToricGeometry/ExistsEtaleCoordsVanishing.lean}
(\srcid{1c87b47bf131dbe99f0958c2abc699184d7a633a}).
The topological group and compactness proofs are respectively
\srcid{ToricGeometry/IsTopologicalGroupComplexPoints.lean}
(\srcid{7e9ca75c5be097e97403b219612289baa50acd6e}) and
\srcid{ToricGeometry/CompactSpacePointsPolAV.lean}
(\srcid{b90f408674114e76130f561572900076bdffdcf6}).
Product holomorphy is
\srcid{ToricGeometry/HolFamilyOfIsPullback.lean}
(\srcid{c343c07f96ec7ec7d2e61d8735cbf9b1b16cf273}).

The quantitative halving limit and extension (CE9) are
\srcid{HomogeneousIdeals/TendstoUniformlyOnDoublingLogOfContraction.lean}
(\srcid{69e6b8b788ca80ded9032668f76eec74e50ed461}).
The local group-law identities, analytic inverse logarithm, discrete
kernel, full span and compact/divisible surjectivity are
\srcid{HomogeneousIdeals/IsLocalHomeomorphAddMonoidHomOfEventuallyEq.lean}
(\srcid{085296514fbf6aaab8595e9593f38e6f4ecc8efc}).
Their actual application to \(P\) is, despite its file name,
\srcid{ToricGeometry/ContinuousOnValue.lean}
(\srcid{967923ce3ac55f13e79463c63e672671f9edcf37}),
the declaration \srcid{AbelianExp.exists\_holUnif}.
The tangent construction is
\srcid{ToricGeometry/TangentBijectionOfValues.lean}
(\srcid{ae10755ee2ab73933551218a57128e702c5cd2c4});
the endomorphism argument is
\srcid{ToricGeometry/ExistsLinearLift.lean}
(\srcid{fb716479e1ebeb2427d806db980225401efb1c21}).
The fixed ordering of these calls is
\srcid{ToricGeometry/ThetaUnifAVOfLaws.lean}
(\srcid{3e496ea337ef9e519bbf43dfe91ef52009f6a960}).
The standard analytic inverse-function and locally uniform limit
theorems are used with the finite-dimensional complex and open-domain
hypotheses displayed above.

\subsection{From the cubical generator to the actual polarising frame}
\label{subsec:ftb-cubical-complex-frame}

Keep the actual exponential \(e:V=\mathbb C^g\to P(\mathbb C)\) and
lattice \(\Lambda\) of
Proposition~\ref{prop:ftb-actual-complex-exponential}.
Let \(L\) be the polarising line bundle, with its prescribed
rigidification at the origin.  The algebraic input in this step is the
source's \(\mathrm{CubeSt}\), applied to this rigidified \(L\).
It supplies a nowhere vanishing algebraic generator \(t\) of
\[
 \Theta(L)=\bigotimes_{I\subset\{0,1,2\}}
              s_I^*L^{\,\sigma_I},\qquad
 s_I(x)=\sum_{a\in I}x_a,\qquad \sigma_I=(-1)^{3-|I|}.        \tag{CF1}
\]
Only this particular bundle and generator are used.

Write \(U_i\) for the algebraic charts of \(L\), and \(g_{ij}\) for their
transition units, in the convention that section coordinates satisfy
\(s_i=g_{ij}s_j\).  Define
\[
 W_i=e^{-1}U_i,\qquad c_{ij}(z)=g_{ij}(e(z)).
\]
For a choice \(j=(j_I)_{I\subset\{0,1,2\}}\), let \(T_j(x)\) be the value
of the \(j\)-coordinate of \(t\) at
\((e(x_0),e(x_1),e(x_2))\).  The charts \(W_i\) cover \(V\).
The \(c_{ij}\) are nonzero holomorphic transition functions and
\[
 c_{ij}c_{jk}=c_{ik},\qquad
 T_j(x)=\prod_I c_{j_Ij'_I}(s_Ix)^{\sigma_I}\,T_{j'}(x).    \tag{CF2}
\]
Each \(T_j\) is holomorphic and nonzero on the open set where all the
\(s_Ix\) lie in \(W_{j_I}\).  Moreover \(W_i,c_{ij}\) are
\(\Lambda\)-periodic, and \(T_j\) is periodic under independent
\(\Lambda\)-translations of the three variables.
These assertions follow by evaluating the algebraic transition and
generator identities: the partial sum of the three points is
\(e(s_Ix)\), and the cube transition unit evaluates to the product in
(CF2).  Thus this is the pulled-back algebraic cube, not a separately
chosen analytic cubical object.

\begin{srlemma}[The source's local third-difference construction]
\label{lem:ftb-cube-local-third-difference}
The preceding actual data admit a radius \(r>0\), nonvanishing
holomorphic functions \(b_i\) on \(W_i\cap B(0,r)\), and a scalar
\(\kappa\ne0\), such that
\[
 b_i=c_{ij}b_j,\qquad
 T_j(x)=\kappa\prod_I b_{j_I}(s_Ix)^{\sigma_I}                \tag{CF3}
\]
for all sufficiently small \(x\) in the relevant cube chart.
\end{srlemma}

\begin{proof}
Choose \(i_0\) with \(0\in W_{i_0}\), and put
\(F(x)=T_{(i_0,\ldots,i_0)}(x)\).
We first give the actual face and doubling argument.
On a face \(x_a=0\), every subset term is paired with the subset
obtained by inserting or deleting \(a\).  The transition products in
(CF2) cancel in pairs.  Hence the corresponding local \(T_j\)'s glue to
a single entire function on that face, periodic under a real-spanning
lattice.  It is constant: choose a real lattice basis, reduce to a
compact fundamental parallelepiped, and apply Liouville's theorem.
Its value at zero is \(F(0)\).  Thus \(F(x)=F(0)\) on each small face.

There is one further glued entire periodic quotient.  Its local
expression is
\[
 \frac{T_{j^+}(2x)
       \prod_I T_{\operatorname{diag}(j_I,j^-_I,j^+_I,i_0)}
                    (s_Ix,s_Ix,-s_Ix)^{\sigma_I}}
      {T_j(x)^3T_{j^-}(-x)}.                               \tag{CF4}
\]
Here the eight charts of \((z,z,-z)\) use, respectively, charts for
\(z,-z,2z,0\): its subset sums are
\(0,z,z,-z,2z,0,0,z\).
In this order, the transition product of this diagonal cube is
\(c(z)^3c(-z)c(2z)^{-1}\).  Substitution in (CF2) cancels every
transition in (CF4).  It is nonzero, entire and lattice-periodic,
hence constant by the same argument.  Its value at zero gives
\[
 F(2x)\prod_I F(s_Ix,s_Ix,-s_Ix)^{\sigma_I}F(0)^3
                         =F(x)^3F(-x).                    \tag{CF5}
\]

Near zero take the analytic logarithm
\(f(x)=\log(F(x)/F(0))\), so \(f(0)=0\).
The face identities give \(f=0\) whenever one of the three vector
variables is zero.  The source obtains
\(f(x)=O(\|x\|^3)\) by applying the one-variable Schwarz estimate
successively to those three variables.  More explicitly, for a fixed
small nonzero \(y\), put \(\rho=R/\|y\|>1\) and consider
\(f(t_0y_0,t_1y_1,t_2y_2)\) on the polydisc
\(|t_a|<\rho\), where it is bounded by \(1\).
Divide first by \(t_0\), then by \(t_1\), then by \(t_2\), using
the face zero and the removable value at zero at each step.
The successive Schwarz estimates give at \(t_0=t_1=t_2=1\)
the bound \(\rho^{-3}=\|y\|^3/R^3\).

Set
\[
 \phi(z)=f(z,z,-z),\qquad
 (\mathcal Dk)(x)=\sum_I\sigma_I k(s_Ix).
\]
Taking logarithms in (CF5) yields, near zero,
\[
       f(2x)+\mathcal D\phi(x)-3f(x)-f(-x)=0.               \tag{CF6}
\]
There is no undetermined logarithmic multiple here: the exponential
of the left side is \(1\), while that side tends to zero.  Shrink so
its modulus is less than \(2\pi\); its possible value
\(2\pi i n\) then has \(n=0\).

Solve
\[
                    k(2z)-3k(z)-k(-z)=-\phi(z).            \tag{CF7}
\]
The source uses a convergent doubling series, not division by formal
Taylor coefficients.  For a cubic-order analytic \(h\), set
\[
 A_0(z)=h(z/2),\quad
 A_{n+1}(z)=3A_n(z/2)+A_n(-z/2),\quad k=\sum_{n\ge0}A_n.
\]
If \(|h(z)|\le C\|z\|^3\) in a ball, then
\(|A_n(z)|\le2^{-n}C\|z\|^3\) there.
The series is locally normally convergent and solves
\(k(2z)-3k(z)-k(-z)=h(z)\) by shifting its summation index.
Use \(h=-\phi\).  Then \(k(0)=0\) and \(k=O(\|z\|^3)\).

Equations (CF6)--(CF7) show that
\(G=f-\mathcal Dk\) satisfies
\(G(2x)=3G(x)+G(-x)\), and \(G=O(\|x\|^3)\).
For \(M(x)=|G(x)|+|G(-x)|\) this gives
\[
 M(x)\le4^n M(2^{-n}x)
               \le2C\|x\|^3\,2^{-n}.
\]
Let \(n\to\infty\).  Therefore \(f=\mathcal Dk\) near zero.
Define
\[
          b_i(z)=c_{ii_0}(z)e^{k(z)},\qquad \kappa=F(0)
\]
on a sufficiently small ball contained in \(W_{i_0}\).
The transition law proves \(b_i=c_{ij}b_j\); exponentiating
\(f=\mathcal Dk\) and applying (CF2) gives (CF3).
All these functions are nonzero, and \(F(0)\ne0\).
\end{proof}

\begin{srlemma}[The actual extension of the frame to all of \(V\)]
\label{lem:ftb-cube-global-frame}
The \(b_i\) extend to nonvanishing holomorphic functions on \(W_i\),
still satisfying their transition law, such that (CF3) holds on every
cube chart, for all \(x\).
\end{srlemma}

\begin{proof}
First extend a local cube identity to the whole region
\[
                    \{x:s_Ix\in B(0,r)\text{ for all }I\}.
\]
This region is convex.  On its cube-chart cover, the quotient of
\(T_j\) by the alternating product of the \(b_{j_I}\) is holomorphic
and agrees on overlaps by (CF2).  Its germ at zero is the constant
\(\kappa\); the identity theorem on this connected region makes the
quotient equal to \(\kappa\) everywhere.

One can now double the radius.  For \(w\in B(0,2r)\) put
\[
                     h(w)=(w/2,w/2,-w/2),\qquad D=\{0,1\}.
\]
The \(D\)-sum is \(w\).  Every other subset sum has norm at most
\(\|w\|/2<r\), and \(\sigma_D=-1\).
On a cube chart \(j\), solve (CF3) for the \(D\)-factor and set
\[
 b'_i(w)=c_{i,j_D}(w)\,
      \frac{\kappa\prod_{I\ne D}
                   b_{j_I}(s_Ih(w))^{\sigma_I}}{T_j(h(w))}. \tag{CF8}
\]
The denominator and every numerator factor are nonzero.
Equations (CF2) and \(c_{ij}c_{jk}=c_{ik}\) show that (CF8) is
independent of the auxiliary cube chart.  Thus these expressions
glue holomorphically; they are nonzero and have the same transition
law.  On \(B(0,r)\), (CF3) proves \(b'_i=b_i\).
The preceding convex-region argument supplies the cube identity
on the larger ball.  Iterate over the balls \(B(0,2^nr)\).
The extensions agree on overlaps, so they glue on all \(W_i\).
For any fixed cube, one sufficiently large ball contains all its
eight subset sums, proving the global form of (CF3).
\end{proof}

\begin{srproposition}[The frame and its exact Appell--Humbert data]
\label{prop:ftb-actual-appell-humbert-frame}
On the same \(e,V,\Lambda\) there are a Hermitian form \(H\), an
integral bilinear form \(E\) on \(\Lambda\), a function
\(\chi:\Lambda\to\mathbb C\) of modulus \(1\), and nonzero holomorphic
frame coordinates \(f_i\) on \(W_i\), with
\[
 \begin{split}
 f_i(z)&=c_{ij}(z)f_j(z),\\
 f_i(z)&=\chi(\lambda)
       \exp\!\left(\pi H(z,\lambda)+\tfrac{\pi}{2}H(\lambda,\lambda)\right)
       f_i(z+\lambda),\\
 E(\lambda,\mu)&=\operatorname{Im}H(\lambda,\mu).             \tag{CF9}
 \end{split}
\]
The \(f_i(0)\) equal the prescribed rigidification coordinates at
zero.  Positivity of \(H\) and perfection of \(E\) are not asserted in
this proposition; they are proved from the polarization and Weil
nondegeneracy in the next step.
\end{srproposition}

\begin{proof}
For \(\lambda\in\Lambda\) define, on any chart containing \(z\),
\[
                    \alpha_\lambda(z)=b_i(z)/b_i(z+\lambda).
\]
Periodicity of \(c_{ij}\) makes this independent of \(i\).
It is an entire nonzero function.  Compare the global cube identity
at \(x=(0,y,z)\) and at \(x+(\lambda,0,0)\).
The cube generator is periodic, and cancellation of its eight
frame factors gives
\[
       \alpha_\lambda(y+z)\alpha_\lambda(0)
                         =\alpha_\lambda(y)\alpha_\lambda(z). \tag{CF10}
\]
An entire multiplicative character \(u\) on the additive vector
space, with \(u(0)=1\), is \(\exp(du_0(z))\).
Indeed differentiation of \(u(z+w)=u(z)u(w)\) gives
\(du_w=u(w)du_0\); the derivative of
\(u(z)\exp(-du_0(z))\) is zero, and its value at zero is \(1\).
Therefore
\[
 \alpha_\lambda(z)=a(\lambda)e^{B_\lambda(z)},\qquad
                       a(\lambda)=\alpha_\lambda(0)\ne0,   \tag{CF11}
\]
where \(B_\lambda\) is complex linear.
The quotient identity
\(\alpha_{\lambda+\mu}(z)=
 \alpha_\lambda(z)\alpha_\mu(z+\lambda)\)
gives, after interchanging \(\lambda,\mu\) if necessary,
\[
 a(\lambda+\mu)=a(\lambda)a(\mu)e^{B_\lambda(\mu)},\qquad
              B_{\lambda+\mu}=B_\lambda+B_\mu.              \tag{CF12}
\]
For the second identity, equality of exponentials of complex linear
forms implies equality of the forms by differentiating at zero.

Here is the explicit normal-form calculation.  Extend the additive
map \(\lambda\mapsto B_\lambda\) real linearly from the full integral
lattice to \(V\); a lattice basis is also a real basis.
The skew relation from (CF12) is
\[
        B_\lambda(\mu)-B_\mu(\lambda)\in2\pi i\mathbb Z.     \tag{CF13}
\]
Its real part vanishes, first on the lattice and then on all of
\(V\times V\) by real linearity.
Splitting a real linear map into its complex linear and conjugate
linear parts therefore writes
\[
                   B_\lambda(z)=A(\lambda,z)+K(z,\lambda),
\]
where \(A\) is symmetric complex bilinear and \(K\) is Hermitian.
Explicitly, the splitting uses the values on the real basis vectors
\(e_j,ie_j\), with coefficients
\((F(e_j)\mp iF(ie_j))/2\); the real skew identity gives the stated
symmetries.  Put \(H=K/\pi\).
Taking imaginary parts in (CF13) shows
\(\operatorname{Im}H(\lambda,\mu)\in\mathbb Z\), and its additivity
in both variables defines \(E\).

The map
\[
 \lambda\longmapsto
 \log|a(\lambda)|-\tfrac12\operatorname{Re}A(\lambda,\lambda)
                 -\tfrac12\operatorname{Re}K(\lambda,\lambda)
\]
is additive by (CF12).  Extend it real linearly to \(V\), and write
this real linear function as \(\operatorname{Re}c(z)\) for a complex
linear \(c\).  Set
\[
 Q(z)=\tfrac12 A(z,z)+c(z),\qquad
 \chi(\lambda)=a(\lambda)
   e^{-A(\lambda,\lambda)/2-K(\lambda,\lambda)/2-c(\lambda)} .
\]
Its modulus is \(1\).  Substituting these definitions into (CF11)
gives
\[
 a(\lambda)e^{B_\lambda(z)}e^{Q(z)-Q(z+\lambda)}
       =\chi(\lambda)e^{\pi H(z,\lambda)+\pi H(\lambda,\lambda)/2}.
\]
Thus \(b_i(z)e^{Q(z)}\) has the required automorphy and transition
law.  Notice \(Q(0)=0\).

Finally choose one \(i_0\) with \(0\in W_{i_0}\), and multiply every
frame coordinate by the single nonzero scalar
\[
 \frac{\text{the prescribed rigidification coordinate in }i_0}
                  {b_{i_0}(0)e^{Q(0)}}.
\]
The algebraic rigidification transition identity makes the value in
every other origin chart correct as well.  This scalar changes
neither the exponential, the lattice, nor the automorphy in (CF9).
\end{proof}

\paragraph{Frozen-source ledger and remaining steps.}
All following files are under
\srcid{src/Development/AbelianVarietyTheory/Proofs/ToricGeometry/}
at \srcid{alpoge-source-7cbf7a0c}.
Actual evaluation of (CF1)--(CF2) is
\srcid{CubeTrivialisationOnUniformisingSpace.lean}
(\srcid{94bdcfed693818f0cdc6f7c89ad5cad13958ca40}).
The periodic face and quotient argument is
\srcid{CubeSectionFacesAndDoubling.lean}
(\srcid{e1fbe94e28161601e128fd1024f1b8bf91f012a2});
it calls the periodic entire-function lemma in
\srcid{../CovolumeBounds/EntireApplyEqOfPeriodicOfSpanReal.lean}.
The series solution to (CF7) is
\srcid{DifferentiableOnDoublingIterate.lean}
(\srcid{04ea97cf72ae242c74021e29294ee9fb81a9f6e3});
the homogeneous cubic estimate is
\srcid{EventuallyEqZeroOfIsBigOCubeOfDoubling.lean}
(\srcid{4b65b8c1681c1d8ee0924fe3afa20faf0cc6a002});
the logarithmic, Schwarz and third-difference assembly is
\srcid{SeriesApplyEqZeroOfFace.lean}
(\srcid{89bfbd686accf0eebbc48ddd5acefd5a5632b35b}).
The local coordinates in (CF3) are constructed in
\srcid{ExistsLocalFrameCompatibleCube.lean}
(\srcid{5c9dc9db6d8d8b04c16fb0643375c36c6dc22444}).
The convex-region and successive radius-doubling construction
(CF8) is \srcid{ExistsGlobalFrameOfLocal.lean}
(\srcid{3fd0ef6e101051ea7a7f5a83620d81c624974784}).
Equations (CF10)--(CF12) are
\srcid{AutomorphyFactorExpLinearOfCube.lean}
(\srcid{9e2f1f482c740af89dc5a110875fb5c07be897fc});
the explicit real-linear splitting and gauge are
\srcid{AppellHumbertNormalFormOfLinearExponent.lean}
(\srcid{1d1f093ab73ab5ad6edb72b9978e2e96148e08f0}).
The actual polarising bundle, same-exponential construction and
origin-rigidification normalization are
\srcid{ExistsAppellHumbertFrameIsRig.lean}
(\srcid{16c73490578c5cc5c0939d1bb71d450477e7be78}).
This proves exactly its Hermitian-frame output.
The later positivity, Weil/perfect-pairing and section-surjectivity
producers are still separate from this chunk.


\subsection{Weil sign, positivity, and all tensor-power sections on the same torus}
\label{subsec:ftb-complex-weil-sections}

Keep exactly the exponential \(e\), lattice \(\Lambda\), frame \(f_i\),
Hermitian form \(H\), unit-modulus function \(\chi\), and integral form
\(E=\operatorname{Im}H|_\Lambda\) constructed in
Proposition~\ref{prop:ftb-actual-appell-humbert-frame}.
Write
\[
 j_\lambda(z)=\chi(\lambda)
       \exp\{\pi H(z,\lambda)+\pi H(\lambda,\lambda)/2\}.
\]
Thus \(f_i(z)=j_\lambda(z)f_i(z+\lambda)\).
The algebraic inputs for this step are precisely
\(\mathrm{Out.WeilSt}\), \srcid{Out.WeilNondegSt}, and
\srcid{Out.GlobalSectionsRankAlgClosed}, in addition to the
multiplication and cube laws already used.
The last law says, for every polarised abelian variety over an
algebraically closed field and every \(n\geq1\),
\(\dim H^0(P,L^n)=n^g\).
We do not replace the source's section comparison by a general
analytic algebraization assertion.

\begin{srlemma}[The actual theta-group scalar and the Weil sign]
\label{lem:ftb-complex-weil-sign}
For \(k\geq1\), \(a,b\in\Lambda\), and
\(x=e(a/k),y=e(b/k)\), the source's algebraic Weil relation is
\[
                 e_{L,k}(x,y)=
                  \exp\{-2\pi i E(a,b)/k\}.                \tag{CW1}
\]
In particular \(E:\Lambda\to\operatorname{Hom}(\Lambda,\mathbb Z)\)
is bijective.
\end{srlemma}

\begin{proof}
Take any algebraic theta-group lift over a test scheme, then restrict
it to a complex point of that test scheme, whose translation is
\(e(a/k)\).  In two charts for \(L\), write \(A_j(z)\) for the value
of its algebraic generator of
\(L_z^k\otimes L_{z+a/k}^{-k}\).  Its transition identity gives
\[
       A_j(z)f_{j_1}(z+a/k)^k
                    =C_a e^{-\pi H(z,a)} f_{j_0}(z)^k      \tag{CW2}
\]
for one scalar \(C_a\) independent of both \(j\) and \(z\).
Here is the source's verification of that scalar.
The quotient
\[
 G_j(z)=
   \frac{A_j(z)f_{j_1}(z+a/k)^k e^{\pi H(z,a)}}
        {f_{j_0}(z)^k}
\]
is holomorphic, and its denominators are nonzero.
The generator and frame transition identities make the \(G_j\)
agree on all overlaps.  Under \(z\mapsto z+\lambda\), their
automorphy factors cancel because
\[
             e^{\pi H(\lambda,a)}=e^{\pi H(a,\lambda)}
\]
for \(\lambda,a\in\Lambda\), by integrality of
\(\operatorname{Im}H(\lambda,a)\).  Hence they glue to a lattice-periodic
entire function and are constant by the fundamental-parallelepiped
and Liouville argument.  Rearranging gives (CW2).

Use (CW2) for the \(a\)-lift at \(z\), the \(b\)-lift at \(z+a/k\),
the \(b\)-lift at \(z\), and the \(a\)-lift at \(z+b/k\).
Cancel the three nonzero intermediate frame values.
The ratio between the two compositions is exactly
\[
 \frac{e^{-\pi H(z,a)}e^{-\pi H(z+a/k,b)}}
      {e^{-\pi H(z,b)}e^{-\pi H(z+b/k,a)}}
       = e^{-2\pi i E(a,b)/k}.                             \tag{CW3}
\]
This fixes the negative sign in (CW1).

This calculation proves the required algebraic generator equality,
not just its value for some chosen lifts.  In the calculation for a
theta-group lift over a test scheme, first restrict that lift to a
complex point of the test scheme; (CW2)--(CW3) are then calculations
on the actual complex variety \(P\).
For the predicate \(\mathrm{WeilIs}\), the source applies
\(\mathrm{WeilSt}\) over the base \(S=\operatorname{Spec}\mathbb C\).
It obtains its lifts and its Weil relation after restriction to an
open \(V\subset S\).  The space on which the resulting regular
generator equality is checked is therefore \(V\times_{\mathbb C}P\).
This space is smooth, reduced and locally of finite type over
\(\mathbb C\).  Formula (CW3), at all its complex points, determines
the scalar in that relation.  Equality of the regular functions
follows on each affine open: their difference lies in every maximal
ideal of a finitely generated complex algebra, hence in its zero
nilradical.  No reducedness assertion about a product with an
arbitrary test scheme is used.
Restriction, pullback, multiplication of generators and constant
units all have exactly their evaluated products in this argument.
This is the source predicate \(\mathrm{WeilIs}\) with scalar (CW1).

To prove perfection, suppose \(k\mid E(a,b)\) for every \(b\in\Lambda\).
Every \(k\)-torsion point is \(e(b/k)\): lift it to \(w\in V\), so
\(kw\in\Lambda\), and take \(b=kw\).
Formula (CW1) therefore says that \(e(a/k)\) pairs trivially with
all \(k\)-torsion points.  The supplied algebraic nondegeneracy
\(\mathrm{WeilNondegSt}\) gives \(e(a/k)=0\).
Thus \(a/k\in\Lambda\), proving
\[
 \bigl(k\mid E(a,b)\text{ for every }b\bigr)
                            \Longrightarrow a\in k\Lambda. \tag{CW4}
\]

For clarity, the source's integral algebra from (CW4) is as follows.
The lattice is finite free over \(\mathbb Z\).
If \(E(a,-)=0\), (CW4) makes every coordinate of \(a\) divisible
by every positive \(k\), so \(a=0\).  Thus \(E\) is injective.
Its cokernel has rank zero, since source and target have equal
finite rank.  For any functional \(\ell\), choose a nonzero integer
\(c\) and \(x\in\Lambda\) with \(E(x,-)=c\ell\).
Then \(E(cx,-)=c^2\ell\), so (CW4) gives \(cx=c^2a'\) for some
\(a'\).  Applying \(E\) and cancelling \(c^2\ne0\) gives
\(E(a',-)=\ell\).  This proves surjectivity.
\end{proof}

For \(n\geq0\), let \(\mathcal T_n\) be the vector space of entire
functions satisfying
\[
                    \theta(z+\lambda)=j_\lambda(z)^n\theta(z).
                                                                    \tag{CW5}
\]
The actual section-to-function map is
\[
            \theta_s(z)=s_i(e(z))/f_i(z)^n
                         \quad\text{whenever }e(z)\in U_i. \tag{CW6}
\]
It is independent of \(i\), by the matching section and frame
transition identities.  It is entire because its local denominator
never vanishes, and (CW5) follows by evaluating at \(z+\lambda\).
The same nonvanishing proves uniqueness of \(\theta_s\).
Conversely, two algebraic sections with the same function have the
same values at every complex point, by surjectivity of \(e\).
The reduced finite-type argument in the proof of
Lemma~\ref{lem:ftb-complex-weil-sign}, applied to each section chart,
therefore makes the two sections equal.  Thus (CW6) is an injective
linear map \(H^0(P,L^n)\to\mathcal T_n\) for every \(n\).

\begin{srlemma}[Positivity from the source's nonzero theta function]
\label{lem:ftb-complex-theta-positivity}
The Hermitian form \(H\) is positive definite.
\end{srlemma}

\begin{proof}
The rank law for \(n=1\) supplies a nonzero algebraic section.
Its function \(\theta\) under (CW6) is nonzero by injectivity.
The continuous function
\[
           |\theta(z)|\exp\{-\pi H(z,z)/2\}
\]
is lattice-periodic: \(|\chi|=1\), and the real part of the
Hermitian automorphy exponent gives the precise quadratic
difference.  It is bounded on a closed fundamental parallelepiped.
Consequently, for some real \(C\geq0\),
\[
                  |\theta(z)|\leq C e^{\pi H(z,z)/2}.       \tag{CW7}
\]
Here \(H(z,z)\) is real.

If \(H(w,w)<0\), then for each fixed \(z\) the entire one-variable
function \(t\mapsto\theta(z+tw)\) has, by (CW7), a bound
\[
 C\exp\{\tfrac\pi2(A+2B|t|-\delta|t|^2)\},
 \qquad \delta=-H(w,w)>0,\quad B=|H(z,w)|.
\]
It tends to zero at infinity.  Liouville's theorem makes it identically
zero.  Varying \(z\) would give \(\theta=0\), a contradiction.
Hence \(H\) is positive semidefinite.

The source next removes its possible nullspace using \(E\).
If \(H(z,z)=0\), positivity of
\(H(z+isv,z+isv)\) for every real \(s\) gives
\[
                 0\leq 2s\,\operatorname{Im}H(z,v)
                                  +s^2H(v,v).
\]
Choose \(s=-c/(d+1)\), where
\(c=\operatorname{Im}H(z,v)\) and \(d=H(v,v)\geq0\).
The resulting expression is
\(-c^2(d+2)/(d+1)^2\), so \(c=0\).
Choose an integral basis of \(\Lambda\), which is a real basis of
\(V\).  In these coordinates the last equality for each lattice
basis vector says that the coordinate row of \(z\) lies in the
left kernel of the matrix of \(E\).  That matrix has nonzero
determinant by Lemma~\ref{lem:ftb-complex-weil-sign}; it remains
invertible over \(\mathbb R\).  Therefore \(z=0\).
This proves positive definiteness.
\end{proof}

\begin{srlemma}[The source's finite Fourier detector]
\label{lem:ftb-complex-theta-fourier-detector}
For \(n\geq1\), every linearly independent family in
\(\mathcal T_n\) has at most \(n^g\) members.
\end{srlemma}

\begin{proof}
We retain the actual symplectic normal form and its gauge.
The full lattice has rank \(2g\).  The alternating perfect form
\(-E\) has an integral symplectic basis \((a_i,b_i)\), so
\(E(a_i,b_j)=-\delta_{ij}\).
The source obtains this basis by the following integral elimination.
Among ideals generated by a value of an alternating form on the
current submodule, choose a maximal one.  A Bezout combination
shows that its generator \(d=B(e,f)\) divides \(B(e,y)\),
\(B(x,f)\), and every value on the common orthogonal complement.
Subtract the corresponding multiples of \(e,f\) to decompose each
vector into their span and that complement.  Induction on rank
gives orthogonal pairs spanning the module.  Their nonzero \(d_i\)
give linear independence; surjectivity of the form forces each
\(d_i\) to be a unit by applying it to a basis coordinate functional.
Rescale the second member of each pair by this unit.  The resulting
basis is symplectic.  This is applied to \(-E\), explaining the
sign of \(E(a_i,b_j)\).

Let \(A\) have columns \(a_i\).  The isotropic relation makes
\(A^tH\overline A=Y\) a real matrix.  It is invertible:
otherwise a real kernel vector of \(Y\), positivity of \(H\),
and the pairing with every \(b_j\) would force that vector to be
zero.  Thus \(A\) is a complex isomorphism.  Write the matrix of
the second half of the basis as \(A\tau\).
The pairing equations give
\[
             Y>0,\qquad Y\,\operatorname{Im}\tau=1,\qquad
             \tau^t=\tau,\qquad \operatorname{Im}\tau>0 .
\]
The symmetry follows by substituting the zero pairings between
the \(b_i\)'s, separately in real and imaginary parts.

For appropriate complex linear \(c\) put
\(Q(z)=\pi z^tYz/2+c(z)\).
Choose \(c(e_i)=\log\chi(a_i)\).  Direct subtraction of \(Q\)
under \(z\mapsto z+e_i\) and \(z\mapsto z+\tau e_j\) shows that
\[
                   F(z)=\theta(Az)e^{-nQ(z)}
\]
is entire, \(1\)-periodic in each coordinate, and
\[
          F(z+\tau e_j)=\beta_j^n
                    e^{-2\pi i n z_j}F(z),\qquad \beta_j\ne0. \tag{CW8}
\]
The map \(\theta\mapsto F\) is injective since \(A\) is invertible
and the gauge exponential never vanishes.

For such an \(F\), define the actual unit-cube Fourier coefficients
\[
             c_m(F)=\int_{[0,1]^g}F(x)e^{-2\pi i m\cdot x}\,dx,
                      \qquad m\in\mathbb Z^g .
\]
The source's shift identity is
\[
 \int_{[0,1]^g}F(x+w)e^{-2\pi i m\cdot x}\,dx
                      =e^{2\pi i m\cdot w}c_m(F).           \tag{CW9}
\]
Indeed the integral of an entire periodic function over the real
unit cube is unchanged by a real translation (Haar translation
on the unit torus).  Differentiation under the integral on this
compact cube makes the integral an entire function of the complex
translation.  Real translation invariance then makes it constant.
Apply this to \(F(z)e^{-2\pi i m\cdot z}\) to get (CW9).

Combining (CW8)--(CW9) gives
\[
       \beta_j^n c_{m+ne_j}(F)
                         =e^{2\pi i m\cdot\tau e_j}c_m(F).  \tag{CW10}
\]
Both multipliers are nonzero.  Vanishing of the coefficients
with \(0\leq m_j<n\) therefore implies vanishing for every
\(m\in\mathbb Z^g\): use integer division of each coordinate
by \(n\) and move by the steps \(\pm ne_j\).
If all coefficients vanish, \(F=0\).  To see this last assertion
on any complex translate of the real torus, (CW9) says that all
its Fourier coefficients vanish.  Fourier uniqueness for continuous
functions on the compact real torus gives zero on that translate,
in particular at its origin.  The translate was arbitrary.

Thus the \(n^g\) coefficients in the residue box define an injective
linear map on this function space into \(\mathbb C^{n^g}\).
Applying this map to a linearly independent family proves the claim.
\end{proof}

\begin{srproposition}[The full fixed-fibre realization]
\label{prop:ftb-full-fixed-fibre-theta-realization}
With the supplied multiplication, cube, Weil, nondegeneracy and
algebraic section-rank laws, the single tuple
\((e,\Lambda,f_i,H,\chi,E)\) above satisfies the full
\(\mathrm{ThetaUnifAV}\) conclusion.  In particular \(H>0\),
\(E\) is perfect, (CW1) has its displayed negative sign, and for
every \(n\geq0\) the rule (CW6) is a bijection between the
algebraic sections of \(L^n\) and \(\mathcal T_n\).
The tangent comparison, linear endomorphism lifts, original
rigidification, and exponential are the ones previously constructed.
\end{srproposition}

\begin{proof}
For \(n\geq1\), the algebraic rank law supplies \(n^g\) independent
sections.  Their functions are independent: a zero linear
combination of functions comes from the same linear combination
of sections by (CW6), so injectivity gives a zero combination of
sections.  Lemma~\ref{lem:ftb-complex-theta-fourier-detector}
then makes these functions span \(\mathcal T_n\).
Otherwise adjoining one function outside their span would give
\(n^g+1\) independent functions.  The corresponding linear combination
of the original algebraic sections gives the required section;
injectivity proves uniqueness.

For \(n=0\), (CW5) says simply that \(\theta\) is lattice-periodic.
It is bounded on a fundamental parallelepiped and hence constant
by Liouville.  The algebraic constant section of \(L^0=\mathcal O_P\)
has exactly this function.  This also gives uniqueness.
The remaining conclusions are precisely the preceding lemmas
and the exponential and frame constructions, whose exponentials
and lattices were unchanged.
\end{proof}

\paragraph{Exact frozen-source ledger.}
Unless otherwise indicated the following paths are under
\srcid{src/Development/AbelianVarietyTheory/Proofs/ToricGeometry/},
at tag \srcid{alpoge-source-7cbf7a0c}.
The scalar comparison (CW2) is
\srcid{ExistsConstOfCover.lean}
(\srcid{352bc50a120c49044d56f2cb6d27897e02e9ded2}),
declaration \srcid{ThetaGroupFrame.thetaGrp\_in\_frame}.
The four-frame cancellation and algebraic Weil relation are
\srcid{WeilIsExpRiemannForm.lean}
(\srcid{1eb1a1affef7a79c6e3ba7d22535430afa9f5930}).
The divisibility step from nondegeneracy is
\srcid{RiemannFormBijectiveOfWeilNondeg.lean}
(\srcid{1f530935aa3e1561f47e7164efbc1e450969a70a});
the finite-free rank and \(c^2\) argument is
\srcid{BilinFormBijectiveOfNondegenerateMod.lean}
(\srcid{996fa526419927f2aa701e51691e6593347aa79e}).
Equality of regular functions from their complex values uses
\srcid{EqZeroOfForallAlgHomComplex.lean}
(\srcid{292d8cb01972525a6c0040b1449efb5090dcca4f}) and
\srcid{IsReducedSchemeOfSmoothOfIsIntegrallyClosed.lean}
(\srcid{14b7a75b14c084dda1e8a2b2b06dd4d26a0c3fba}).

Construction and injectivity of (CW6) are
\srcid{SectionsPowThetaFunctions.lean}
(\srcid{54dfefb784a86c4772781b543aa65b04d961f540}).
The nonzero rank-one section and positivity assembly are
\srcid{ExistsIndepSections.lean}
(\srcid{fa66eca4089a8ad5a3bf75a329926da7b4bc58d8}) and
\srcid{IsAHOfPolarisation.lean}
(\srcid{22ea6df6cd8176d9647db8c52e101b63203fd5a2});
the periodic metric bound, negative-direction argument and
nullspace removal are
\srcid{ReFormPosOfThetaNeZero.lean}
(\srcid{bda58e9df3981e4de3d25e29bb8531c5f4f93e1e}).

The integral symplectic elimination is
\srcid{../CovolumeBounds/DualRestrictSurjectiveOfIsTorsionFreeQuotient.lean}
(\srcid{369a733eead2b44d17430d13ba108b43be4e8f81}).
Its application, the exact \(-E\) sign, period matrix and gauge are
\srcid{IsUnitOfLagrangian.lean}
(\srcid{b2457d0adf49b8651e52f5bf0a955de91cbdefda}),
declaration \srcid{SiegelNormalForm.exists\_normal\_form}.
The complex Fourier shift is
\srcid{FourierCoeffUnitCubeTranslateOfEntirePeriodic.lean}
(\srcid{90e515c3c1229348b46c587391faf4f2f1dfbedd});
Fourier uniqueness on translates is
\srcid{EntirePeriodicEqZeroOfFourierCoeffEqZero.lean}
(\srcid{ba105521538d00b051b34f19676d77e0f374ca90}).
The residue-box argument is
\srcid{CardLePowOfLinearIndependent.lean}
(\srcid{20085cbe9a9a42f684df9924f67d1e8fcc10a49f});
the normal-form transport is
\srcid{CardLePowOfLinearIndependentThetaPrincipal.lean}
(\srcid{fe36db1b4f2c9d4f78040330571ea9d7ca87ef58}).
Surjectivity including \(n=0\) is
\srcid{MemSpanOfCardLe.lean}
(\srcid{f21dc5ff7a5efc6072c172914c28130acaafe764}).
Final same-tuple assembly is
\srcid{ThetaUnifAVOfParts.lean}
(\srcid{d260f9fb5b92e8635a1893372bb0ad47bdacb5c6}),
called in the fixed order by
\srcid{ThetaUnifAVOfLaws.lean}
(\srcid{3e496ea337ef9e519bbf43dfe91ef52009f6a960}).
This completes that producer from its five displayed algebraic laws.
The construction of families of period matrices and the reverse
construction of an algebraic variety from a given period matrix
are separate later producers.
